\chapter{Analitik AI}
\section{Deteksi multi-model}
Setiap stream dapat menjalankan beberapa model YOLO bersamaan. Per model:

\noindent\begin{tabularx}{\linewidth}{@{}l X@{}}
\toprule
Parameter & Arti \\ \midrule
Kunci & Nama pendek unik; hitungan memakai \code{kunci:kelas}. \\
Ambang keyakinan & Deteksi di bawah nilai ini diabaikan. \\
Tiap N frame & Model dijalankan setiap N frame; di antaranya hasil terakhir dipakai. \\
Filter kelas & Hanya kelas terpilih yang ditampilkan dan dihitung. \\
Warna & Warna kotak di video beranotasi. \\
\bottomrule
\end{tabularx}

\section{Hitungan grup dan prompt}
Berkas \code{infra/\allowbreak{}worker/\allowbreak{}prompts.\allowbreak{}toml} dimuat ulang otomatis saat disimpan.
\begin{lstlisting}
[count_groups]
orang = ["pedestrian", "people", "person"]
kendaraan = ["car", "van", "truck", "bus", "motor"]

[prompts]
crowd = "{orang} people and {kendaraan} vehicles found. What are the people doing?"

[moondream]
active = ["caption", "crowd"]     # bergiliran tiap siklus
interval_sec = 1.0

[translate]
model = "llama3.2:3b"             # caption diterjemahkan ke Bahasa Indonesia
\end{lstlisting}
Entri \code{"kelas"} mencocokkan kelas itu di semua model; \code{"kunci:kelas"} hanya di satu model. Placeholder prompt: nama grup, \code{\{counts\}}, \code{\{time\}}.

\section{Caption dan tanya-jawab visual}
Caption Moondream hanya berjalan untuk stream yang sedang dibuka di mode fokus (hemat GPU). Tanya-jawab visual memakai \code{GCS\_VLM\_MODEL}: \code{moondream} (cepat, Inggris, diterjemahkan bila \code{GCS\_\allowbreak{}TRANSLATE\_\allowbreak{}MODEL} diisi) atau VLM multibahasa seperti \code{qwen2.5vl:3b} yang menjawab langsung dalam Bahasa Indonesia.

\section{Speech-to-text}
Mikrofon di panel tanya-jawab memakai faster-whisper di server (model \code{small}, CPU). Bila tidak terpasang, konsol memakai pengenal suara browser.
